\section*{Bab \ref{ch:g8:the-oceans} --- Samudra}

\begin{solution}{exo:g8:the-oceans:1}
Mula-mula \omterm{def:g8:the-oceans:seafloor}{paparan benua}, lalu \omterm{def:g8:the-oceans:seafloor}{lereng benua}, kemudian \omterm{def:g8:the-oceans:seafloor}{dataran abisal}, dan akhirnya \omterm{def:g8:the-oceans:seafloor}{punggung samudra}.
\end{solution}

\begin{solution}{exo:g8:the-oceans:2}
(a) Salah: sekitar \qty{3700}{m}, yaitu \qty{3.7}{km}. (b) Benar. (c)
Salah: penguapan mengambil airnya dan meninggalkan garamnya, sehingga
airnya menjadi lebih asin. (d) Benar.
\end{solution}

\begin{solution}{exo:g8:the-oceans:3}
(a) \omterm{def:g8:the-oceans:seafloor}{Punggung samudra}. (b) \omterm{def:g8:the-oceans:salinity}{Salinitas}. (c) \omterm{def:g8:the-oceans:current}{Pilin samudra}. (d) \omterm{def:g8:the-oceans:seafloor}{Palung samudra}.
\end{solution}

\begin{solution}{exo:g8:the-oceans:4}
(a) Berlawanan arah jarum jam. (b) Hangat: Arus Teluk; dingin: Arus
Kanari. (c) \omterm{def:g8:winds-global-circulation:global-winds}{Angin pasat}.
\end{solution}

\begin{solution}{exo:g8:the-oceans:5}
(a) $35 \div 1000 \times 100 = 3.5~\%$. (b) $20 \times 35 = \qty{700}{g}$.
(c) $1000 \div 35 \approx 28.6$: sekitar \qty{29}{kg}.
\end{solution}

\begin{solution}{exo:g8:the-oceans:6}
(a) $1500 \times 6 = \qty{9000}{m}$. (b) Bunyi turun sampai ke dasar lalu
naik kembali: bunyi menempuh kedalaman itu dua kali. (c) $9000 \div 2 =
\qty{4500}{m}$. (d) \omterm{def:g8:the-oceans:seafloor}{Dataran abisal} (palung akan lebih dalam, punggung
atau paparan lebih dangkal).
\end{solution}

\begin{solution}{exo:g8:the-oceans:7}
(a) Antara sekitar \qty{3500}{m} dan \qty{5000}{m}. (b) Sekitar
\qty{4400}{km} sepanjang profil, dengan kedalaman sekitar \qty{2300}{m}.
(c) Sekitar $4500 - 2300 = \qty{2200}{m}$ (antara 1200 dan
\qty{2700}{m}, bergantung pada \omterm{def:g1:landscapes-around-us:landform}{dataran} yang dipilih). (d) \omterm{def:g2:water-on-earth:ocean}{Samudra} hanya
beberapa kilometer dalamnya, tetapi ribuan kilometer lebarnya: jika
digambar dengan skala yang sama, seluruh profil akan lebih tipis daripada
sebuah garis.
\end{solution}

\begin{solution}{exo:g8:the-oceans:8}
\omterm{def:g2:water-on-earth:ocean}{Laut} A lebih asin: Matahari dan udara yang kering menyebabkan
\emph{penguapan} yang kuat, airnya pergi, \emph{garam tertinggal}, dan
hampir tidak ada \omterm{def:g4:clouds-rain-snow:precipitation}{hujan} atau \omterm{def:g1:landscapes-around-us:water}{sungai} yang membawa \emph{\omterm{def:g2:water-on-earth:salt-fresh}{air tawar}} kembali.
\omterm{def:g2:water-on-earth:ocean}{Laut} B kurang asin: \omterm{def:g1:landscapes-around-us:water}{sungai-sungainya} dan \omterm{def:g4:clouds-rain-snow:precipitation}{hujannya} membawa banyak
\emph{\omterm{def:g2:water-on-earth:salt-fresh}{air tawar}}, yang mengencerkan garamnya.
\end{solution}

\begin{solution}{exo:g8:the-oceans:9}
(a) Arus Teluk, ke utara di sepanjang pantai, lalu ke timur laut. (b) Ke
timur (timur laut), bersama Arus Atlantik Utara. (c) Di sepanjang pantai
Afrika barat laut (serta Spanyol dan Portugal), bersama Arus Kanari.
(d) Dingin.
\end{solution}

\begin{solution}{exo:g8:the-oceans:10}
(a) $6.4 \times 24 = 153.6$: sekitar \qty{154}{km}. (b) $153.6 \times 7
\approx \qty{1075}{km}$. (c) Arus Teluk sedikit lebih cepat daripada
orang yang berjalan kaki, tetapi arus itu tidak pernah berhenti, siang
dan malam.
\end{solution}

\begin{solution}{exo:g8:the-oceans:11}
(a) \qty{3.7}{km} dan \qty{11}{km}. (b) $11000 \div 3700 \approx 3.0$
kali. (c) Bunyi menempuh $2 \times 11000 = \qty{22000}{m}$; $22000
\div 1500 \approx 14.7$: sekitar \qty{15}{s}. (d) $1000 = 10^3$ tahun.
\end{solution}

\begin{solution}{exo:g8:the-oceans:12}
(a) Di dekat \omterm{def:g3:continents-and-oceans:equator}{khatulistiwa}, \omterm{def:g8:winds-global-circulation:global-winds}{angin pasat} mendorong air tropis yang hangat ke
barat. Di tepi barat Atlantik, air itu berbelok ke utara bersama Arus
Teluk, sebuah arus hangat, di sepanjang Florida dan Amerika Utara. \omterm{def:g8:winds-global-circulation:global-winds}{Angin
baratan} lalu mendorongnya menyeberangi Atlantik ke arah timur laut,
sampai ke \omterm{def:g2:water-on-earth:ocean}{laut-laut} di lepas pantai Norwegia. (b) Di sana air itu
melepaskan panasnya ke udara yang dingin: air itu mendingin, dan menjadi
lebih asin ketika es \omterm{def:g2:water-on-earth:ocean}{laut} terbentuk, karena garamnya tertinggal di dalam
air. Air yang dingin dan asin itu rapat, sehingga tenggelam ke dasar.
(c) Air itu merayap ke selatan di sepanjang dasar \omterm{def:g2:water-on-earth:ocean}{samudra}, ke arah
Antarktika, lalu masuk ke \omterm{def:g2:water-on-earth:ocean}{Samudra} Hindia dan \omterm{def:g2:water-on-earth:ocean}{Samudra} Pasifik, tempat air
itu perlahan-lahan naik kembali.
\end{solution}

\begin{solution}{pb:g8:the-oceans:1}
\textbf{1.} Di \omterm{def:g2:water-on-earth:ocean}{Samudra} Atlantik, di \omterm{def:g3:continents-and-oceans:equator}{belahan bumi utara}.

\textbf{2.} Searah jarum jam; pelampung itu mula-mula terbawa ke utara.

\textbf{3.} Tahap A (Arus Teluk) hangat: arusnya datang dari \omterm{def:g5:weather-and-climate:tropics}{daerah
tropis}. Tahap C (Arus Kanari) dingin: arusnya datang dari \omterm{def:g6:latitude-longitude:latitude}{lintang} yang
lebih tinggi, menuju \omterm{def:g3:continents-and-oceans:equator}{khatulistiwa}.

\textbf{4.} Tahap B: \omterm{def:g8:winds-global-circulation:global-winds}{angin baratan}. Tahap D: \omterm{def:g8:winds-global-circulation:global-winds}{angin pasat}.

\textbf{5.} $3800 \div 100 = 38$ hari.

\textbf{6.} B: $2100 \div 20 = 105$ hari. C: $3100 \div 20 = 155$ hari.
D: $6100 \div 20 = 305$ hari.

\textbf{7.} $3800 + 2100 + 3100 + 6100 = \qty{15100}{km}$.

\textbf{8.} Arus Teluk adalah arus yang sempit dan cepat di sisi barat
\omterm{def:g8:the-oceans:current}{pilin}; arus-arus di sisi lainnya lebar dan lambat.

\textbf{9.} $38 + 105 + 155 + 305 = 603$ hari.

\textbf{10.} $603 \div 30 = 20.1$: sekitar \textbf{20 bulan}.

\textbf{11.} Setelah tahap A, B, dan C, sudah berlalu
$38 + 105 + 155 = 298$ hari; hari ke-365 jatuh pada tahap D (Arus
\omterm{def:g3:continents-and-oceans:equator}{Khatulistiwa} Utara), $365 -
298 = 67$ hari setelah tahap itu dimulai.

\textbf{12.} $603 - 365 = 238$ hari lagi, yaitu $238 \div 30 \approx 8$
bulan lagi.
\end{solution}
